% !PRINTED 271
\weberpage{271}
To attain this goal it suffices here too that, when $n$ has no square divisor, $\Phi_{a,0,\partial}$, and when $n$ is a prime, $\Phi_{1,0,n}$ remain finite for $q=0$.

In these calculations we make use of the expansions:
\weq{6}{\begin{aligned}
uv &= q^{-\frac{n+1}{24}}\Pi(1+q^{n(2h+1)})(1+q^{2h+1}),\\
u_1v_1 &= q^{-\frac{n+1}{24}}\Pi(1-q^{n(2h+1)})(1-q^{2h+1}),\\
u_2v_2 &= 2q^{\frac{n+1}{12}}\Pi(1+q^{2hn})(1+q^{2h}),
\end{aligned}}
from which, by expansion in powers of $q$:
\weq{7}{\begin{aligned}
uv &= q^{-\frac{n+1}{24}}(1+q+q^3+q^4+q^5+q^6+q^7+2q^8\dots),\\
u_1v_1 &= q^{-\frac{n+1}{24}}(1-q-q^3+q^4-q^5+q^6-q^7+2q^8\dots),\\
u_2v_2 &= 2q^{\frac{n+1}{12}}(1+q^2+q^4+2q^6+2q^8+\cdots).
\end{aligned}}
The latter formulae are correct for $n>7$ (for $n=3,5,7$ the terms from $q^3, q^5, q^7$ on are to be modified). For $n=7$, $n=23$ these expressions show that $A$ itself remains finite for $q=0$, from which the equations follow for these cases:
\weq{8}{\begin{aligned}
n&=7, & A&=0,\\
n&=23, & A&=1.
\end{aligned}}
By a simple calculation one finds further still:
\weq{9}{\begin{aligned}
n&=31, & (A^2-B)^2-A&=0,\\
n&=47, & A^2-A-B&=2,\\
n&=71, & A^3-4A^2+2A-B&=1.
\end{aligned}}
The calculation turns out not quite so simple for a composite $n$. Thus for $n=15$, for example, one must take account of the condition of finiteness for $q=0$ not only for $\Phi_{1,0,15}$, but also for $\Phi_{3,0,5}$. These two, however, suffice. One obtains in this way:
\weq{10}{n=15, \quad A^3-AB+1=0.}
2. If $n \equiv 3$ (mod 8), the same conclusions are to be drawn, if we set:
\weq{11}{\begin{aligned}
4A &= u^2v^2 - u_1^2v_1^2 - u_2^2v_2^2,\\
B &= u^2v^2u_1^2v_1^2 + u^2v^2u_2^2v_2^2 - u_1^2v_1^2u_2^2v_2^2,
\end{aligned}}
%CONT>>
